% GENERATED FILE. Edit canonical lessons, then run npm run generate:books.
% canonical-source-hash: fnv1a64:7cfc4cef4d8ac627
% canonical-lessons: ML-C142-kalivu, ML-C142-sakti, ML-C142-avakasam, ML-C142-avasaram, ML-C142-upakaranam

\chapter{Ability, Strength, Right, Opportunity, Tool}
\label{ch:ml-ability-strength-right-opportunity-tool}
\hlchaptermodality{142}

\begin{chapteropening}
\textbf{By the end of this chapter:} \emph{I can say ability, strength, a right, an opportunity and a tool.}

Complete the last lesson of chapter 142: I can say ability, strength, a right, an opportunity and a tool.
\end{chapteropening}

\section[\texorpdfstring{\ml{കഴിവ്} (kaḻivŭ)}{കഴിവ് (kaḻivŭ)}]{\ml{കഴിവ്} (kaḻivŭ) — ability, skill}

\label{lesson:ML-C142-kalivu}



\begin{quote}
Before the new one: say the Malayalam for Eid, a feast day, then the Malayalam for a festival.

Say \textquotedblleft{}weather.\textquotedblright{} (\emph{Kālāvastha} — \emph{kālam}, \textquotedblleft{}time,\textquotedblright{} plus a word for a state.)
\end{quote}

\subsection*{You'll want to know: \ml{കഴിവ്}}
\textbf{\ml{കഴിവ്}} — \emph{kaḻivŭ} — \textquotedblleft{}ability, skill\textquotedblright{}.

\begin{cousinweb}
\textbf{\ml{കഴിവ്}} (\emph{kaḻivŭ}) is native Dravidian, the noun of the verb \textquotedblleft{}to be able\textquotedblright{}; DEDR 1357 lists the root with Tamil \emph{kaḻi}.
\end{cousinweb}

\begin{grammarlens}[title={What you've built}]
One more word for this chapter.
\end{grammarlens}

\subsection*{Your turn}
\begin{itemize}
\raggedright
  \item \emph{Say it:} \emph{kaḻivŭ}
  \item \emph{Say it:} \emph{kaḻivŭ}, once more
  \item \emph{Say it:} say \emph{kaḻivŭ}
  \item \emph{Recall:} say the Malayalam for Eid, a feast day, then the Malayalam for a festival, then say \emph{kaḻivŭ} again
\end{itemize}

\begin{tcolorbox}[breakable,colback=teal!4,colframe=teal!35!black,title={Before you move on}]
What does \ml{കഴിവ്} mean? (\textquotedblleft{}Ability, skill\textquotedblright{}.) Say it once more.
\end{tcolorbox}
\section[\texorpdfstring{\ml{ശക്തി} (śakti)}{ശക്തി (śakti)}]{\ml{ശക്തി} (śakti) — strength, power}

\label{lesson:ML-C142-sakti}



\begin{quote}
Before the new one: say the Malayalam for a festival, then the Malayalam for ability.
\end{quote}

\subsection*{You'll want to know: \ml{ശക്തി}}
\textbf{\ml{ശക്തി}} — \emph{śakti} — \textquotedblleft{}strength, power\textquotedblright{}.

\begin{cousinweb}
\textbf{\ml{ശക്തി}} (\emph{śakti}) $\leftarrow$ Sanskrit \emph{śakti}, \textquotedblleft{}power\textquotedblright{}, taken over whole as a tatsama (CDIAL 12250).
\end{cousinweb}

\begin{grammarlens}[title={What you've built}]
One more word for this chapter.
\end{grammarlens}

\subsection*{Your turn}
\begin{itemize}
\raggedright
  \item \emph{Say it:} \emph{śakti}
  \item \emph{Say it:} \emph{śakti}, once more
  \item \emph{Say it:} say \emph{śakti}
  \item \emph{Recall:} say the Malayalam for a festival, then the Malayalam for ability, then say \emph{śakti} again
\end{itemize}

\begin{tcolorbox}[breakable,colback=teal!4,colframe=teal!35!black,title={Before you move on}]
What does \ml{ശക്തി} mean? (\textquotedblleft{}Strength, power\textquotedblright{}.) Say it once more.
\end{tcolorbox}
\section[\texorpdfstring{\ml{അവകാശം} (avakāśaṁ)}{അവകാശം (avakāśaṁ)}]{\ml{അവകാശം} (avakāśaṁ) — a right}

\label{lesson:ML-C142-avakasam}



\begin{quote}
Before the new one: say the Malayalam for ability, then the Malayalam for strength.

Which part of \emph{kālāvastha}, \textquotedblleft{}weather,\textquotedblright{} means \textquotedblleft{}time\textquotedblright{}? (\emph{Kālam}.)
\end{quote}

\subsection*{You'll want to know: \ml{അവകാശം}}
\textbf{\ml{അവകാശം}} — \emph{avakāśaṁ} — \textquotedblleft{}a right\textquotedblright{}.

\begin{cousinweb}
\textbf{\ml{അവകാശം}} (\emph{avakāśaṁ}) $\leftarrow$ Sanskrit \emph{avakāśa}, \textquotedblleft{}space, room, opportunity\textquotedblright{}, taken over whole as a tatsama, with Malayalam's noun ending \emph{-ṁ} (CDIAL 728). Room to act became a right.
\end{cousinweb}

\begin{grammarlens}[title={What you've built}]
One more word for this chapter.
\end{grammarlens}

\subsection*{Your turn}
\begin{itemize}
\raggedright
  \item \emph{Say it:} \emph{avakāśaṁ}
  \item \emph{Say it:} \emph{avakāśaṁ}, once more
  \item \emph{Say it:} say \emph{avakāśaṁ}
  \item \emph{Recall:} say the Malayalam for ability, then the Malayalam for strength, then say \emph{avakāśaṁ} again
\end{itemize}

\begin{tcolorbox}[breakable,colback=teal!4,colframe=teal!35!black,title={Before you move on}]
What does \ml{അവകാശം} mean? (\textquotedblleft{}A right\textquotedblright{}.) Say it once more.
\end{tcolorbox}
\section[\texorpdfstring{\ml{അവസരം} (avasaraṁ)}{അവസരം (avasaraṁ)}]{\ml{അവസരം} (avasaraṁ) — an opportunity, a chance}

\label{lesson:ML-C142-avasaram}



\begin{quote}
Before the new one: say the Malayalam for strength, then the Malayalam for a right.
\end{quote}

\subsection*{You'll want to know: \ml{അവസരം}}
\textbf{\ml{അവസരം}} — \emph{avasaraṁ} — \textquotedblleft{}an opportunity, a chance\textquotedblright{}.

\begin{cousinweb}
\textbf{\ml{അവസരം}} (\emph{avasaraṁ}) $\leftarrow$ Sanskrit \emph{avasara}, \textquotedblleft{}occasion, opportunity\textquotedblright{}, taken over whole as a tatsama, with Malayalam's noun ending \emph{-ṁ} (CDIAL 858).
\end{cousinweb}

\begin{grammarlens}[title={What you've built}]
One more word for this chapter.
\end{grammarlens}

\subsection*{Your turn}
\begin{itemize}
\raggedright
  \item \emph{Say it:} \emph{avasaraṁ}
  \item \emph{Say it:} \emph{avasaraṁ}, once more
  \item \emph{Say it:} say \emph{avasaraṁ}
  \item \emph{Recall:} say the Malayalam for strength, then the Malayalam for a right, then say \emph{avasaraṁ} again
\end{itemize}

\begin{tcolorbox}[breakable,colback=teal!4,colframe=teal!35!black,title={Before you move on}]
What does \ml{അവസരം} mean? (\textquotedblleft{}An opportunity, a chance\textquotedblright{}.) Say it once more.
\end{tcolorbox}
\section[\texorpdfstring{\ml{ഉപകരണം} (upakaraṇaṁ)}{ഉപകരണം (upakaraṇaṁ)}]{\ml{ഉപകരണം} (upakaraṇaṁ) — a tool}

\label{lesson:ML-C142-upakaranam}



\begin{quote}
Before the new one: say the Malayalam for a right, then the Malayalam for an opportunity.

Besides \emph{divasaṁ}, which more formal word for \textquotedblleft{}day\textquotedblright{} does Malayalam keep? (\emph{Dinam}.)
\end{quote}

\subsection*{You'll want to know: \ml{ഉപകരണം}}
\textbf{\ml{ഉപകരണം}} — \emph{upakaraṇaṁ} — \textquotedblleft{}a tool\textquotedblright{}.

\begin{cousinweb}
\textbf{\ml{ഉപകരണം}} (\emph{upakaraṇaṁ}) $\leftarrow$ Sanskrit \emph{upakaraṇa}, \textquotedblleft{}implement\textquotedblright{}, taken over whole as a tatsama, with Malayalam's noun ending \emph{-ṁ} (CDIAL 2138).
\end{cousinweb}

\begin{grammarlens}[title={What you've built}]
One more word for this chapter.
\end{grammarlens}

\subsection*{Your turn}
\begin{itemize}
\raggedright
  \item \emph{Say it:} \emph{upakaraṇaṁ}
  \item \emph{Say it:} \emph{upakaraṇaṁ}, once more
  \item \emph{Say it:} say \emph{upakaraṇaṁ}
  \item \emph{Recall:} say the Malayalam for a right, then the Malayalam for an opportunity, then say \emph{upakaraṇaṁ} again
\end{itemize}

\begin{tcolorbox}[breakable,colback=teal!4,colframe=teal!35!black,title={Before you move on}]
What does \ml{ഉപകരണം} mean? (\textquotedblleft{}A tool\textquotedblright{}.) Say it once more.
\end{tcolorbox}
